\documentclass[12pt]{article}

\usepackage[margin=1.2in]{geometry}
\usepackage{amsmath,amsthm,amssymb}
\usepackage{microtype}
\usepackage[authoryear,round]{natbib}
\usepackage{xcolor}
\definecolor{linkgreen}{RGB}{0,120,100}
\usepackage[colorlinks=true,linkcolor=linkgreen,citecolor=linkgreen,urlcolor=linkgreen]{hyperref}

% Reduce title font from \LARGE (default) to \Large so title fits in 2 lines
\makeatletter
\renewcommand{\maketitle}{%
  \begin{center}%
    {\Large\bfseries \@title \par}%
    \vskip 1.5em%
    {\normalsize \@author \par}%
    \vskip 0.8em%
    {\normalsize \@date}%
  \end{center}%
  \vskip 2em%
}
\makeatother

% Theorem environments
\theoremstyle{plain}
\newtheorem{proposition}{Proposition}
\newtheorem{lemma}{Lemma}
\newtheorem{corollary}{Corollary}
\theoremstyle{definition}
\newtheorem{assumption}{Assumption}
\newtheorem{definition}{Definition}
\theoremstyle{remark}
\newtheorem{remark}{Remark}

\title{Capacity-Constrained Dominant Firm Pricing\\
in Segmented Markets: A Theoretical Framework}
\author{%
  \texttt{pAI-Econ-claude} (\texttt{theoretical-economics-claude-skill})\\[4pt]
  \small\url{https://github.com/maxwell2732/pAI-Econ-claude}\\[2pt]
  \small Claude Sonnet 4.6%
}
\date{Draft: June 15, 2026}

\begin{document}

\maketitle

\section{Introduction}

In many consumer goods markets in developing economies, a single dominant
branded firm competes against a fragmented competitive fringe that differs in
strength across geographic regions.  A stylized version of this structure arises
when one firm holds a substantially larger market share than any individual rival
and the remaining supply comes from many small, price-taking producers; urban and
rural markets are physically separated so that direct consumer arbitrage is
prohibitively costly.  When the dominant firm cannot fully serve all markets
simultaneously due to binding production or distribution capacity, a natural
question arises: how does the asymmetry in regional competitive conditions shape
the firm's pricing strategy across markets, and does intensified competition in
one region spill over to affect prices in the other?

This paper shows that the answer is yes, and that the mechanism operates entirely
through the shadow price of the dominant firm's capacity constraint.  We study a
static model of third-degree price discrimination by a capacity-constrained
dominant firm facing a competitive fringe in two geographically segmented
markets.  The dominant firm sets prices in each market, anticipating the fringe
firms' competitive supply response.  A binding capacity constraint links the two
markets: using more capacity in one market necessarily reduces the amount
available for the other.  This linkage generates a shadow cost that appears in
both markets' first-order conditions, transmitting competitive shocks across
market boundaries even when direct arbitrage between markets is prohibitively
costly.

The standard dominant firm model, developed for single-market settings, predicts
no cross-market pricing effects when markets are geographically segmented: the
firm's pricing decision in each market depends only on local demand and
competitive conditions.  Similarly, \citet{Varian1985}'s canonical analysis of
third-degree price discrimination focuses on a monopolist with no fringe
competition and does not address how one market's equilibrium responds to
changes in the competitive structure of another.  Neither framework can generate
the cross-market spillovers we study, because neither introduces the binding
capacity constraint that provides the linking mechanism.

Our main results are threefold.  First, we characterize the unique
dominant firm equilibrium in closed form (Proposition~\ref{prop:existence}), and
show that equilibrium prices in each market depend on both local demand and fringe
parameters \emph{and} on a common capacity shadow price that summarizes market
scarcity.  Second, we derive the price structure under segmentation.  When the
urban market has a higher residual demand intercept per unit of fringe competition
(Assumption~A4$'$), the dominant firm charges a strictly higher price in the urban
market.  Remarkably, this price gap is independent of the capacity level and of
the dominant firm's cost---it is determined entirely by the asymmetry in fringe
supply structures (Proposition~\ref{prop:pricegap}).  Third, and most
importantly, we show that an increase in rural fringe competition reduces the
dominant firm's prices in \emph{both} markets, while simultaneously widening the
urban-rural price gap (Proposition~\ref{prop:crossmarket}).  The urban price
falls because the dominant firm reallocates capacity toward the urban market where
margins are higher, reducing the shadow price of capacity.  The price gap widens
because only the rural price absorbs the direct effect of stronger rural
competition.  We also show that capacity expansion by the dominant firm reduces
both prices uniformly, leaving the price gap unchanged
(Proposition~\ref{prop:capacity}).

On welfare, we show that market segmentation---relative to a uniform pricing
benchmark---reduces total market output and total welfare when fringe competition
is structurally stronger in the rural market (Proposition~\ref{prop:welfare}).
The welfare loss arises through the fringe: because the dominant firm produces
its full capacity $K$ in both regimes, the welfare difference is driven entirely
by fringe output.  Stronger rural fringe competition means rural output is more
sensitive to the rural price; the output loss from the dominant firm's supra-uniform
urban price exceeds the output gain from its sub-uniform rural price.  Applying
\citeauthor{Varian1985}'s criterion to total market output delivers the welfare
ranking.  The distributional picture is nuanced: rural consumers benefit from the
lower prices the dominant firm charges in the rural market, while urban consumers
pay higher prices.  A regulator who cares about spatial equity faces a direct
trade-off between rural consumer welfare and aggregate efficiency.

Our final result characterizes the boundary between the constrained and
unconstrained regimes (Proposition~\ref{prop:boundary}).  Price levels adjust
continuously as the capacity constraint becomes slack.  However, the cross-market
pricing sensitivity---the degree to which urban prices respond to rural
competition---is \emph{discontinuous} at the capacity threshold $\bar{K}$: it
converges to $c/(2B)$ (where $B$ is the sum of all fringe and demand slope
parameters) as the constraint approaches slack, then drops to zero the moment
the constraint is lifted.  This means a barely-constrained dominant firm and an
unconstrained dominant firm behave qualitatively differently in how they respond
to competitive shocks in distant markets.

This paper contributes to two strands of the industrial organization literature.
The first is the literature on dominant firm competition and price leadership,
initiated by \citet{Gaskins1971}, who analyzed how a dominant firm optimally
prices to slow fringe entry over time.  We study a static version of this
framework, and show that introducing a binding capacity constraint across two
segmented markets generates new structural predictions absent from the
single-market model.  The second strand is \citeauthor{Varian1985}'s
(\citeyear{Varian1985}) analysis of third-degree price discrimination, which
established that welfare rises if and only if total output rises under
discrimination.  We apply this criterion in a dominant firm setting where total
output is the sum of dominant firm and fringe output, and show that the welfare
direction is determined by fringe output differences across regimes---a
mechanism that does not appear in Varian's monopoly framework.

The paper proceeds as follows.  Section~\ref{sec:model} presents the model.
Section~\ref{sec:equilibrium} characterizes the equilibrium and the price gap.
Section~\ref{sec:comparative} derives the comparative statics.
Section~\ref{sec:welfare} analyzes welfare.  Section~\ref{sec:discussion}
discusses robustness and extensions.  Section~\ref{sec:conclusion} concludes.
Proofs are in the Appendix.

%----------------------------------------------------------------------
\section{Model}
\label{sec:model}
%----------------------------------------------------------------------

\subsection{Environment}

A single dominant firm $D$ sells a homogeneous good in two geographically
segmented markets: urban ($u$) and rural ($r$).  Each market $m \in \{u,r\}$
also contains a competitive fringe of atomistic price-taking firms.  The timing
is as follows.  First, $D$ commits to a price pair $(p_u, p_r) \in \mathbb{R}_+^2$.
Second, fringe firms in each market $m$ observe $p_m$ and supply $S_m(p_m)$.
Third, residual demand in market $m$ is $R_m(p_m) = D_m(p_m) - S_m(p_m)$, and
$D$ supplies $q_m = R_m(p_m)$ in market $m$.  The game is static and complete
information.

\textbf{Demand.}  Consumer demand in market $m$ is linear:
\begin{equation}
  D_m(p_m) = a_m - b_m p_m, \quad a_m,\, b_m > 0.
\end{equation}
We denote the choke price $\bar{P}_m = a_m/b_m$.

\textbf{Fringe supply.}  Each fringe firm $i$ in market $m$ has cost
$c_i^F(q_i) = q_i^2/(2g_m)$, yielding individual supply $g_m p_m$ at price $p_m$.
Aggregating over a continuum of measure $N_m \to \infty$ with $e_m \equiv N_m g_m$
held fixed, aggregate fringe supply is:
\begin{equation}
  S_m(p_m) = e_m p_m, \quad e_m > 0.
\end{equation}
The fringe supply elasticity parameter $e_m$ is the key structural variable
characterizing competitive conditions in market $m$.

\textbf{Residual demand.}  The dominant firm's residual demand in market $m$ is:
\begin{equation}
  R_m(p_m) = D_m(p_m) - S_m(p_m) = a_m - (b_m + e_m)p_m.
\end{equation}

\textbf{Dominant firm.}  The dominant firm has constant marginal cost $c > 0$
and a binding production and distribution capacity constraint:
\begin{equation}
  q_u + q_r \equiv R_u(p_u) + R_r(p_r) \leq K.
  \label{eq:capacity}
\end{equation}
The dominant firm maximizes profit $\Pi_D = (p_u - c)R_u(p_u) + (p_r - c)R_r(p_r)$
subject to~\eqref{eq:capacity} and $R_m(p_m) \geq 0$.

Let $B \equiv b_u + e_u + b_r + e_r$ denote the sum of all slope and fringe
parameters, and define the \emph{unconstrained capacity threshold}:
\begin{equation}
  \bar{K} \equiv \frac{a_u + a_r - Bc}{2}.
  \label{eq:Kbar}
\end{equation}
At $K = \bar{K}$, the dominant firm's capacity exactly accommodates its
unconstrained optimal quantities in both markets; for $K < \bar{K}$ the
constraint binds.  Note that $\bar{K} < (a_u + a_r)/2$ since $Bc > 0$,
so Assumption~\ref{ass:A3} below also implies $K < (a_u + a_r)/2$.

\subsection{Assumptions}

\begin{assumption}[Market Segmentation]
\label{ass:A1}
Per-unit arbitrage between markets costs $\tau > 0$, with $\tau > p_u^* - p_r^*$
in equilibrium.  No consumer arbitrage occurs.
\end{assumption}

\begin{assumption}[Dominant Firm Cost Advantage]
\label{ass:A2}
$c < p_m^*$ for $m \in \{u, r\}$, so the dominant firm earns positive profit
in both markets.
\end{assumption}

\begin{assumption}[Binding Capacity Constraint]
\label{ass:A3}
$K < \bar{K}$, where $\bar{K}$ is defined in~\eqref{eq:Kbar}.  The capacity
constraint binds in equilibrium and $\lambda^* > 0$.
\end{assumption}

\begin{assumption}[Interior Solution]
\label{ass:H2}
$a_m > (b_m + e_m)(c + \lambda^*)$ for $m \in \{u, r\}$, so $q_m^* > 0$:
the dominant firm supplies a strictly positive quantity in both markets.
\end{assumption}

\begin{assumption}[Asymmetric Residual Demand Intercepts]
\label{ass:A4prime}
$\dfrac{a_u}{b_u + e_u} > \dfrac{a_r}{b_r + e_r}$: the urban market's residual
demand intercept exceeds the rural market's.
\end{assumption}

\begin{remark}
Under symmetric consumer demand ($a_u = a_r$, $b_u = b_r$), Assumption~\ref{ass:A4prime}
reduces to $e_r > e_u$: the dominant firm charges more in the market with
the less elastic fringe.  Asymmetric demand can reverse this ordering (see
Section~\ref{sec:discussion}), which is why the general condition on residual
demand intercepts is the correct sufficient condition.
\end{remark}

\subsection{Equilibrium Concept}

\begin{definition}[Dominant Firm Equilibrium]
A \emph{dominant firm equilibrium} is a price pair $(p_u^*, p_r^*)$ such that:
(i) each fringe firm in market $m$ supplies $s_m^* = e_m p_m^*$ (best response
to $p_m^*$ for a price-taker with quadratic cost); (ii) $(p_u^*, p_r^*)$ solves
$\max \Pi_D$ subject to~\eqref{eq:capacity} and $R_m(p_m) \geq 0$; (iii)
markets clear: $q_m^* + s_m^* = D_m(p_m^*)$.
\end{definition}

This is the price-leadership equilibrium of the two-stage game in which $D$
chooses prices first and fringe firms supply second.  It coincides with the
subgame-perfect Nash equilibrium of that game: in stage~2, fringe firms'
dominant strategy is to supply at the announced price; in stage~1, $D$
maximizes profit given the anticipated fringe response.

\subsection{Social Planner Benchmark}

For the welfare analysis of Section~\ref{sec:welfare}, the relevant benchmark is
the \emph{uniform-price equilibrium}: $D$ is constrained to charge the same price
$p$ in both markets.  The uniform-price optimum $p^{U*}$ solves the same
constrained profit-maximization problem with $p_u = p_r = p$.  Under the binding
capacity constraint:
\begin{equation}
  p^{U*} = \frac{c + \lambda^{U*}}{2} + \frac{a_u + a_r}{2B},
  \quad \lambda^{U*} = \frac{a_u + a_r - 2K}{B} - c > 0.
  \label{eq:uniform}
\end{equation}
Note that $p^{U*} = (c + \lambda^{U*})/2 + (a_u + a_r)/(2B)$, which lies between
$p_r^*$ and $p_u^*$ in the segmented equilibrium (shown in Proposition~\ref{prop:welfare}).

%----------------------------------------------------------------------
\section{Equilibrium and Price Structure}
\label{sec:equilibrium}
%----------------------------------------------------------------------

\subsection{Existence and Characterization}

The Lagrangian for the dominant firm's problem is:
\begin{equation}
  \mathcal{L}(p_u, p_r, \lambda) = \sum_{m \in \{u,r\}} (p_m - c - \lambda)R_m(p_m) + \lambda K.
\end{equation}
Since $R_m(p_m) = a_m - (b_m + e_m)p_m$, the profit function is strictly
concave in $(p_u, p_r)$ for each $\lambda$: the Hessian is diagonal with entries
$-2(b_m + e_m) < 0$.  The first-order conditions for $m \in \{u, r\}$ are:
\begin{equation}
  R_m(p_m^*) + (p_m^* - c - \lambda^*)\,R_m'(p_m^*) = 0,
\end{equation}
where $R_m'(p) = -(b_m + e_m)$.  Solving explicitly:
\begin{equation}
  p_m^* = \frac{c + \lambda^*}{2} + \frac{a_m}{2(b_m + e_m)}.
  \label{eq:pm}
\end{equation}
The shadow price $\lambda^*$ is determined by the binding capacity constraint
$R_u(p_u^*) + R_r(p_r^*) = K$:
\begin{equation}
  \lambda^* = \frac{a_u + a_r - 2K}{B} - c.
  \label{eq:lambda}
\end{equation}
Assumption~\ref{ass:A3} ensures $\lambda^* > 0$.

\begin{proposition}[Existence and Uniqueness]
\label{prop:existence}
Under Assumptions~\ref{ass:A1}--\ref{ass:H2}, there exists a unique dominant firm
equilibrium.  Letting $B \equiv b_u + e_u + b_r + e_r$, the equilibrium prices,
shadow price, and quantities are:
\begin{align}
  p_m^* &= \frac{c + \lambda^*}{2} + \frac{a_m}{2(b_m + e_m)}, \quad m \in \{u,r\}, \\
  \lambda^* &= \frac{a_u + a_r - 2K}{B} - c \;>\; 0, \\
  q_m^* &= \frac{a_m - (b_m + e_m)(c + \lambda^*)}{2} \;>\; 0, \quad
  s_m^* = e_m p_m^* > 0.
\end{align}
\end{proposition}

\begin{proof}
Existence follows from the Weierstrass theorem applied to $\Pi_D$ over the compact
feasible set defined by~\eqref{eq:capacity} and $R_m \geq 0$.  Uniqueness follows
from strict concavity: the Hessian is $\mathrm{diag}(-2(b_u+e_u),\,-2(b_r+e_r)) \prec 0$.
KKT conditions are necessary and sufficient.  Equation~\eqref{eq:lambda} gives
$\lambda^* > 0$ from Assumption~\ref{ass:A3}; $q_m^* > 0$ from Assumption~\ref{ass:H2}.
\end{proof}

The equilibrium price in each market has two components: a common term
$(c + \lambda^*)/2$ that reflects the firm's effective marginal cost (production
cost $c$ plus the shadow price of capacity $\lambda^*$), and a market-specific
term $a_m/[2(b_m + e_m)]$ that reflects the depth of local demand net of fringe
competition.  The shadow price $\lambda^*$ is the mechanism linking the two
markets: it rises when capacity is scarcer (smaller $K$) and falls as the
constraint relaxes.

\subsection{The Urban-Rural Price Gap}

\begin{proposition}[Urban-Rural Price Gap]
\label{prop:pricegap}
In the equilibrium of Proposition~\ref{prop:existence}, the urban-rural price
gap is:
\begin{equation}
  p_u^* - p_r^* = \frac{a_u}{2(b_u + e_u)} - \frac{a_r}{2(b_r + e_r)}.
  \label{eq:gap}
\end{equation}
Under Assumption~\ref{ass:A4prime}, $p_u^* > p_r^*$.  Under symmetric demand
($a_u = a_r = a$, $b_u = b_r = b$) and $e_r > e_u$:
\begin{equation}
  p_u^* - p_r^* = \frac{a}{2}\left[\frac{1}{b+e_u} - \frac{1}{b+e_r}\right] > 0.
  \label{eq:gap-symmetric}
\end{equation}
The price gap is independent of the capacity level $K$, the shadow price
$\lambda^*$, and the dominant firm's marginal cost $c$.
\end{proposition}

\begin{proof}
Subtract the formula for $p_r^*$ from $p_u^*$ in equation~\eqref{eq:pm}.
The common term $(c + \lambda^*)/2$ cancels exactly, leaving~\eqref{eq:gap}.
The sign under Assumption~\ref{ass:A4prime} is immediate.  Independence from
$K$ and $c$ follows because neither appears in~\eqref{eq:gap}.
\end{proof}

The price gap is determined entirely by the structure of residual demand in
each market: specifically, the ratio of the demand intercept to the combined
slope and fringe elasticity.  The capacity shadow price $\lambda^*$ shifts
both prices upward by the same amount, so it cancels in the difference.
This means the dominant firm's level of capacity scarcity has no effect on
how much more it charges in the urban market than in the rural market.

\begin{remark}
\label{rem:gap-asymmetry}
Under asymmetric demand, Assumption~\ref{ass:A4prime} cannot be replaced by
the simpler condition $e_r > e_u$ alone.  If the rural market is sufficiently
larger ($a_r \gg a_u$), then $a_r/(b_r + e_r) > a_u/(b_u + e_u)$ even when
$e_r > e_u$, and $p_r^* > p_u^*$.  Assumption~\ref{ass:A4prime} is the general
sufficient condition.
\end{remark}

%----------------------------------------------------------------------
\section{Comparative Statics}
\label{sec:comparative}
%----------------------------------------------------------------------

\subsection{Cross-Market Effects of Fringe Competition}

The paper's main result concerns how an increase in rural fringe competition---a
rise in $e_r$---affects equilibrium in both markets.

\begin{proposition}[Cross-Market Comparative Statics]
\label{prop:crossmarket}
In the equilibrium of Proposition~\ref{prop:existence}, under
Assumption~\ref{ass:A3} ($K < \bar{K}$):
\begin{align}
  \frac{d\lambda^*}{de_r} &= -\frac{a_u + a_r - 2K}{B^2} < 0, \label{eq:dlambda}\\[4pt]
  \frac{dp_u^*}{de_r} &= \frac{1}{2}\frac{d\lambda^*}{de_r} < 0, \label{eq:dpu}\\[4pt]
  \frac{dp_r^*}{de_r} &= \frac{1}{2}\frac{d\lambda^*}{de_r} - \frac{a_r}{2(b_r + e_r)^2} < 0,
    \label{eq:dpr}\\[4pt]
  \frac{d(p_u^* - p_r^*)}{de_r} &= \frac{a_r}{2(b_r + e_r)^2} > 0, \label{eq:dgap}\\[4pt]
  \frac{dq_u^*}{de_r} &= \frac{(b_u+e_u)(a_u + a_r - 2K)}{2B^2} > 0. \label{eq:dqu}
\end{align}
Furthermore, if $K \geq \bar{K}$ (constraint non-binding), then $dp_u^*/de_r = 0$
and $dq_u^*/de_r = 0$: all cross-market effects in
\eqref{eq:dlambda}--\eqref{eq:dqu} are generated exclusively by the binding
capacity constraint.
\end{proposition}

\begin{proof}
See Appendix~\ref{app:proofs}, Proof of Proposition~\ref{prop:crossmarket}.
\end{proof}

The mechanism behind Proposition~\ref{prop:crossmarket} involves two distinct
channels.  The \emph{shadow price channel} (equation~\eqref{eq:dpu}): when $e_r$
rises, the dominant firm's rural residual demand shrinks, making rural capacity less
valuable.  The capacity shadow price $\lambda^*$ falls, reducing the effective
marginal cost and pulling both prices down by the same amount
$(1/2)(d\lambda^*/de_r) < 0$.  The \emph{direct channel}
(the second term in equation~\eqref{eq:dpr}): the increase in $e_r$ also directly
flattens rural residual demand, reducing the rural price by an additional
$a_r/[2(b_r + e_r)^2]$.  This direct channel is present only in the rural market,
which is why the rural price falls more than the urban price and the gap widens
(equation~\eqref{eq:dgap}).  Meanwhile, as rural becomes less profitable, the
dominant firm shifts its limited capacity toward the urban market ($q_u^*$ rises,
$q_r^*$ falls), reinforcing the urban price reduction through higher urban supply.

The last part of the proposition---that the cross-market effect is exactly zero
when $K \geq \bar{K}$---is the key identification result.  Two dominant firms
that are observationally similar in their price levels may behave qualitatively
differently: the constrained firm transmits rural competitive shocks to urban
prices; the unconstrained firm does not.

\subsection{Capacity Comparative Statics and Gap Invariance}

\begin{proposition}[Capacity Comparative Statics]
\label{prop:capacity}
In the equilibrium of Proposition~\ref{prop:existence}, under
Assumption~\ref{ass:A3}:
\begin{equation}
  \frac{dp_m^*}{dK} = -\frac{1}{B} < 0, \quad m \in \{u, r\}.
  \label{eq:dpdk}
\end{equation}
Consequently, $d(p_u^* - p_r^*)/dK = 0$: the urban-rural price gap is invariant
to the dominant firm's capacity level.  As $K \to \bar{K}^-$:
$p_m^* \to p_m^{nc} \equiv c/2 + a_m/[2(b_m + e_m)]$, the unconstrained
dominant firm price.
\end{proposition}

\begin{proof}
From~\eqref{eq:lambda}, $d\lambda^*/dK = -2/B$.  From~\eqref{eq:pm},
$dp_m^*/dK = (1/2)(d\lambda^*/dK) = -1/B$.  Gap invariance follows because
the price gap~\eqref{eq:gap} contains no $K$.  The limit $p_m^* \to p_m^{nc}$
follows from $\lambda^* \to 0$ as $K \to \bar{K}^-$.
\end{proof}

Proposition~\ref{prop:capacity} has a direct policy implication: capacity
investment by the dominant firm (improved distribution infrastructure, new
production facilities) reduces prices uniformly across markets, but leaves the
urban-rural price differential unchanged.  The capacity shadow price acts as a
common cost reduction that both markets share equally.  By contrast,
Proposition~\ref{prop:crossmarket} shows that changes in fringe supply structures
\emph{do} affect the gap.  These are thus distinct policy instruments: capacity
subsidies improve the \emph{level} of prices, while policies that strengthen rural
fringe competition (reducing barriers to local entry) affect the price
\emph{structure}.

%----------------------------------------------------------------------
\section{Welfare Analysis}
\label{sec:welfare}
%----------------------------------------------------------------------

\subsection{Total Output and Welfare}

Let $Q \equiv q_u + q_r + s_u + s_r$ denote total market output (dominant firm
plus fringe, across both markets).  Under the binding capacity constraint,
dominant firm output equals $K$ in both the segmented and uniform-price regimes.
Welfare differences therefore originate entirely in fringe output:
$Q^{\mathrm{SEG}} - Q^{\mathrm{UNI}} = [e_u p_u^* + e_r p_r^*] - [(e_u + e_r)p^{U*}]
= e_u(p_u^* - p^{U*}) + e_r(p_r^* - p^{U*})$.

\begin{proposition}[Welfare: Segmentation Reduces Output and Welfare]
\label{prop:welfare}
Under Assumptions~\ref{ass:A1}--\ref{ass:A4prime}, symmetric consumer demand
($a_u = a_r = a$, $b_u = b_r = b$), and with the capacity constraint binding in
both regimes:
\begin{equation}
  Q^{\mathrm{SEG}} - Q^{\mathrm{UNI}} = \frac{a(e_r - e_u)}{2B}
  \left[\frac{e_u}{b + e_u} - \frac{e_r}{b + e_r}\right] < 0.
  \label{eq:Qdiff}
\end{equation}
Consequently, $W^{\mathrm{SEG}} < W^{\mathrm{UNI}}$: market segmentation reduces
total welfare relative to uniform pricing.

In terms of distributional welfare effects: rural consumers gain
($CS_r^{\mathrm{SEG}} > CS_r^{\mathrm{UNI}}$, since $p_r^* < p^{U*}$); urban
consumers lose ($CS_u^{\mathrm{SEG}} < CS_u^{\mathrm{UNI}}$, since
$p_u^* > p^{U*}$).
\end{proposition}

\begin{proof}
See Appendix~\ref{app:proofs}, Proof of Proposition~\ref{prop:welfare}.
\end{proof}

The welfare result is an application of \citeauthor{Varian1985}'s
(\citeyear{Varian1985}) criterion to the dominant firm setting: total welfare
increases with total output when all units in $[Q^{\mathrm{SEG}},\, Q^{\mathrm{UNI}}]$
carry positive net social value (which holds here since each ``lost'' unit has
demand price above fringe marginal cost at the equilibrium).  The content of
Proposition~\ref{prop:welfare} relative to \citeauthor{Varian1985} is the
identification of the fringe output differential as the welfare driver: in Varian's
monopoly framework, total output equals the monopolist's choice; here, dominant
firm output is $K$ in both regimes, so welfare is entirely determined by the
fringe's response to price differences.  The sign of the fringe output differential
follows from Assumption~\ref{ass:A4prime}: stronger rural competition means rural
fringe supply is more price-sensitive, so the output loss from the high urban price
exceeds the output gain from the low rural price.

The distributional finding---rural consumers benefit from segmentation while
aggregate welfare falls---creates a policy tension.  A regulator imposing uniform
pricing improves aggregate welfare, but at the cost of raising rural prices.  This
tension is present whenever price discrimination by a dominant firm directs
lower prices toward lower-income consumers in the structurally more competitive
market.

\subsection{Boundary Behavior}

\begin{proposition}[Capacity Boundary]
\label{prop:boundary}
As $K \to \bar{K}^-$, equilibrium prices converge continuously to the
unconstrained dominant firm prices: $\lim_{K \to \bar{K}^-} p_m^* = p_m^{nc}$
and $\lambda^* \to 0$.  However, the cross-market comparative static is
discontinuous at $K = \bar{K}$:
\begin{equation}
  \lim_{K \to \bar{K}^-} \frac{dp_u^*}{de_r} = -\frac{c}{2B} \;\neq\; 0
  \;=\; \frac{dp_u^*}{de_r}\bigg|_{K \geq \bar{K}}.
  \label{eq:boundary}
\end{equation}
The magnitude of the jump is $c/(2B)$, which vanishes as $c \to 0$.
\end{proposition}

\begin{proof}
Price continuity follows from $\lambda^*(K) \to 0$ continuously.  For the
comparative static: at $K < \bar{K}$,
$dp_u^*/de_r = -(a_u + a_r - 2K)/(2B^2)$.  As $K \to \bar{K}^-$, $a_u + a_r - 2K \to Bc$
(using $\bar{K} = (a_u + a_r - Bc)/2$), giving the limit $-c/(2B)$.  At
$K \geq \bar{K}$, the constraint is slack, $\lambda^* = 0$, and $p_u^* = c/2 +
a_u/[2(b_u + e_u)]$ is independent of $e_r$, so $dp_u^*/de_r = 0$.
\end{proof}

Proposition~\ref{prop:boundary} shows that the capacity constraint functions as
a structural switch rather than a continuous dial.  A dominant firm with capacity
just below $\bar{K}$ still transmits rural competition shocks to urban prices
(with magnitude approaching $c/(2B)$), while a firm with capacity just above
$\bar{K}$ transmits nothing.  This qualitative discontinuity has empirical content:
firms in periods of capacity constraint should exhibit cross-market pricing
sensitivity, while the same firms in slack capacity periods should not.

%----------------------------------------------------------------------
\section{Discussion}
\label{sec:discussion}
%----------------------------------------------------------------------

\subsection{Robustness of the Main Results}

\textbf{Asymmetric demand.}  Propositions~\ref{prop:crossmarket}
and~\ref{prop:capacity} (cross-market comparative statics and capacity gap
invariance) are robust to asymmetric consumer demand ($a_u \neq a_r$, $b_u \neq b_r$).
The price gap formula~\eqref{eq:gap} holds generally; the symmetric demand
specialization is used only to obtain a clean closed form for
equation~\eqref{eq:gap-symmetric} and for the welfare quantity difference
in~\eqref{eq:Qdiff}.  The welfare direction under asymmetric demand depends
on whether Assumption~\ref{ass:A4prime} is satisfied.  When $a_r/(b_r + e_r) > a_u/(b_u + e_u)$
(i.e., Assumption~\ref{ass:A4prime} fails), the dominant firm charges more in
the rural market, the fringe supplies more in rural at the higher price, and
total output under segmentation can exceed the uniform-price output---reversing
the welfare direction.  Assumption~\ref{ass:A4prime} is thus the correct
general condition for the welfare loss result, not the simpler $e_r > e_u$.

\textbf{Convex dominant firm cost.}  Replace $C(q) = cq$ with
$C(q) = cq + (\gamma/2)q^2$.  Under a binding capacity constraint, total
dominant firm output is $K$ regardless of the cost function's convexity.
The effective marginal cost at the constraint becomes $c + \gamma K$, which
is a constant at the optimum.  Proposition~\ref{prop:crossmarket} therefore
holds with $c$ replaced by $c + \gamma K$: cross-market effects survive
convex costs, and Proposition~\ref{ass:A3} (the capacity constraint needed
to generate them) remains the operative assumption.

\textbf{Multiple markets.}  With $N > 2$ markets, the equilibrium is
$p_m^* = (c + \lambda^*)/2 + a_m/[2(b_m + e_m)]$ with $\lambda^* = (\sum_m a_m - 2K)/B - c$
and $B = \sum_m (b_m + e_m)$.  The pairwise price gap
$p_m^* - p_{m'}^*$ contains no $K$, so gap invariance extends immediately to
$N$ markets.  The cross-market comparative static $dp_m^*/de_{m'} = d\lambda^*/(2de_{m'}) < 0$
holds for any market pair as long as the capacity constraint binds.

\subsection{Extensions}

\textbf{Strategic capacity investment.}  The model treats $K$ as exogenous.
An important extension is to allow the dominant firm to choose $K$ in a prior
stage.  By Proposition~\ref{prop:capacity}, capacity investment reduces prices
uniformly without affecting the price gap; if the dominant firm internalizes the
gap structure, it may strategically choose $K$ to maintain $K < \bar{K}$,
preserving the cross-market mechanism.  Alternatively, a Dixit-type
\citep{Gaskins1971} argument suggests the dominant firm may under-invest in
capacity to keep $\lambda^*$ high and maintain the urban price premium.

\textbf{Dynamic fringe entry.}  The fringe supply elasticity $e_m$ is fixed in
this model.  A dynamic extension in the spirit of \citet{Gaskins1971} would
allow $e_m$ to evolve endogenously as fringe firms enter in response to price
levels.  The cross-market mechanism identified here would then generate dynamic
spillovers: a dominant firm threatened by rural entry would reduce urban prices
as a by-product of its rural competitive response, even without any urban entry
threat.

\textbf{Empirical identification.}  Proposition~\ref{prop:crossmarket} provides
a testable prediction: the cross-market price sensitivity $dp_u^*/de_r$ is
strictly negative when the capacity constraint binds ($K < \bar{K}$) and exactly
zero when it does not.  Panel data on regional prices combined with capacity
utilization data could test this prediction.  The capacity constraint boundary
$\bar{K}$ depends on demand and cost parameters; seasonal demand variation
in consumer goods markets could provide natural variation in whether the
constraint binds, offering a quasi-experimental test of the mechanism.

\subsection{Limitations}

Several features of real markets are absent from this model.

\textbf{Market structure.}  The model abstracts from strategic interaction
between large firms.  Many real markets that might appear to fit the
dominant-firm-plus-fringe description are better characterized as oligopolies.
China's national carbonated beverage market is a case in point: Coca-Cola held
approximately 61\% and PepsiCo approximately 30\% of market share in 2024,
making PepsiCo a strategic rival rather than an atomistic price-taker.  The
dominant firm equilibrium concept is appropriate only when a single firm is
genuinely dominant and all remaining supply comes from price-taking producers.
Applying the model's predictions to markets with a small number of large
strategic rivals would require an oligopoly framework.

The model further assumes a static, complete-information game; dynamic
considerations (brand loyalty, reputation, repeated price competition) are
not captured.  The welfare analysis
uses the representative consumer framework within each market and thus cannot
address within-market distributional effects.  The assumption of a single
homogeneous product precludes product differentiation across markets, which may
be important for branded beverages.  Finally, the model takes the dominant firm's
cost advantage (Assumption~\ref{ass:A2}) as exogenous; a richer model would
derive cost advantage from economies of scale or sunk investments, which would
interact with the capacity constraint.

%----------------------------------------------------------------------
\section{Conclusion}
\label{sec:conclusion}
%----------------------------------------------------------------------

We have analyzed how a capacity-constrained dominant firm sets prices across two
geographically segmented markets with asymmetric fringe competition.  Our main
results are as follows.  The equilibrium price gap between markets is determined
entirely by the structure of residual demand---not by the firm's capacity level
or cost (Proposition~\ref{prop:pricegap}).  An increase in rural fringe
competition reduces prices in both markets while widening the price gap, through
the shadow price of the shared capacity constraint (Proposition~\ref{prop:crossmarket}).
Capacity expansion reduces both prices uniformly but leaves the gap unchanged
(Proposition~\ref{prop:capacity}).  Market segmentation reduces total welfare,
with rural consumers benefiting and urban consumers losing
(Proposition~\ref{prop:welfare}).  Finally, the cross-market pricing mechanism
switches on and off discontinuously at the capacity threshold
(Proposition~\ref{prop:boundary}).

The central message is that a binding capacity constraint converts what would
otherwise be two independent pricing problems into a single linked system.  The
shadow price of capacity is the mechanism of linkage, and its presence or absence
is a qualitative structural feature of the market.  For competition policy, the
implication is that evaluating a dominant firm's pricing behavior in one market
requires knowing whether the firm faces binding capacity constraints---a condition
that standard single-market analysis ignores.

Natural directions for future work include a dynamic model in which the dominant
firm invests strategically in capacity, fringe firms respond to prices through
entry and exit, and the capacity shadow price evolves over time.  An empirical
companion to this paper would use variation in capacity utilization (e.g., seasonal
demand shocks) to identify the cross-market mechanism predicted by
Proposition~\ref{prop:crossmarket}.

%----------------------------------------------------------------------
\appendix
\section{Proofs}
\label{app:proofs}
%----------------------------------------------------------------------

\textbf{Proof of Proposition~\ref{prop:crossmarket}.}

\emph{Equations~\eqref{eq:dlambda}--\eqref{eq:dpr}.}  Differentiate
$\lambda^* = (a_u + a_r - 2K)/B - c$ with respect to $e_r$.  Since
$\partial B / \partial e_r = 1$:
\begin{equation*}
  \frac{d\lambda^*}{de_r} = -\frac{a_u + a_r - 2K}{B^2}.
\end{equation*}
Under Assumption~\ref{ass:A3}, $K < \bar{K} = (a_u + a_r - Bc)/2$, so
$a_u + a_r - 2K > a_u + a_r - 2\bar{K} = Bc > 0$.  Thus
$d\lambda^*/de_r < 0$, proving~\eqref{eq:dlambda}.

For~\eqref{eq:dpu}: differentiate $p_u^* = (c + \lambda^*)/2 + a_u/[2(b_u + e_u)]$
with respect to $e_r$.  Since $a_u$, $b_u$, $e_u$ do not depend on $e_r$:
\begin{equation*}
  \frac{dp_u^*}{de_r} = \frac{1}{2}\frac{d\lambda^*}{de_r} < 0.
\end{equation*}
For~\eqref{eq:dpr}: differentiate $p_r^* = (c + \lambda^*)/2 + a_r/[2(b_r + e_r)]$
with respect to $e_r$:
\begin{equation*}
  \frac{dp_r^*}{de_r} = \frac{1}{2}\frac{d\lambda^*}{de_r} - \frac{a_r}{2(b_r + e_r)^2} < 0.
\end{equation*}
Both terms are negative, so $dp_r^*/de_r < 0$.

\emph{Equation~\eqref{eq:dgap}.}  Subtract:
$d(p_u^* - p_r^*)/de_r = dp_u^*/de_r - dp_r^*/de_r = a_r/[2(b_r + e_r)^2] > 0$.

\emph{Equation~\eqref{eq:dqu}.}  Since $q_u^* = [a_u - (b_u + e_u)(c + \lambda^*)]/2$:
\begin{equation*}
  \frac{dq_u^*}{de_r} = -\frac{b_u + e_u}{2}\frac{d\lambda^*}{de_r}
  = \frac{(b_u + e_u)(a_u + a_r - 2K)}{2B^2} > 0.
\end{equation*}

\emph{Non-binding case.}  When $K \geq \bar{K}$, the constraint is slack and
$\lambda^* = 0$, so $p_u^* = c/2 + a_u/[2(b_u + e_u)]$, which is independent
of $e_r$.  Thus $dp_u^*/de_r = 0$ and $dq_u^*/de_r = 0$. \hfill$\square$

\bigskip

\textbf{Proof of Proposition~\ref{prop:welfare}.}

\emph{Step~1: price comparisons with uniform benchmark.}  From~\eqref{eq:pm}
and~\eqref{eq:uniform}, under symmetric demand ($a_u = a_r = a$, $b_u = b_r = b$):
\begin{equation*}
  p_u^* - p^{U*} = \frac{a}{2(b+e_u)} - \frac{a}{2} \cdot \frac{2}{2b + e_u + e_r}
  = \frac{a}{2}\left[\frac{1}{b+e_u} - \frac{2}{2b+e_u+e_r}\right].
\end{equation*}
Since $e_r > e_u$ (from Assumption~\ref{ass:A4prime} under symmetric demand):
$1/(b+e_u) > 2/(2b+e_u+e_r)$, so $p_u^* > p^{U*}$.  Similarly
$p_r^* < p^{U*}$.  Thus $p^{U*}$ lies strictly between $p_r^*$ and $p_u^*$.

\emph{Step~2: fringe output differential.}  Under the binding constraint in
both regimes, $q_u^* + q_r^* = q^{U*} = K$.  Total output differs only through
fringe: $Q^{\mathrm{SEG}} - Q^{\mathrm{UNI}} = e_u(p_u^* - p^{U*}) + e_r(p_r^* - p^{U*})$.

\emph{Step~3: sign.}  Substituting the closed forms (algebra omitted for
brevity) yields~\eqref{eq:Qdiff}.  Under Assumption~\ref{ass:A4prime} ($e_r > e_u$
under symmetric demand), $e_u/(b+e_u) < e_r/(b+e_r)$, so the bracket is negative,
giving $Q^{\mathrm{SEG}} < Q^{\mathrm{UNI}}$.

\emph{Step~4: welfare ranking via Varian criterion.}  Total welfare is
$W = \int_0^Q P(x)\,dx - cK - \sum_m C_m^F(s_m^*)$,
where $P(x)$ is the inverse aggregate demand.  Consider any unit $x$ in the range
$(Q^{\mathrm{SEG}},\, Q^{\mathrm{UNI}})$ that is produced under uniform pricing but
not under segmentation.  This unit has aggregate demand price $P(x) > 0$.
Its social cost is the fringe marginal cost $s_m/e_m$ at that output level, which
at equilibrium equals $p_m^*$.  Since $P(x) \geq p_m^*$ for $x$ in this range
(units below equilibrium quantity have demand price above equilibrium price), each
lost unit carries positive net social value.  Hence $W^{\mathrm{SEG}} < W^{\mathrm{UNI}}$.

\emph{Distributional effects.}  $CS_m^{\mathrm{SEG}} - CS_m^{\mathrm{UNI}}
= -(D_m(p_m^*) + D_m(p^{U*}))(p_m^* - p^{U*})/2$ approximately.
Since $p_u^* > p^{U*}$, urban consumer surplus falls; since $p_r^* < p^{U*}$,
rural consumer surplus rises. \hfill$\square$

%----------------------------------------------------------------------
\begin{thebibliography}{99}

\bibitem[Gaskins(1971)]{Gaskins1971}
Gaskins, D.W.\ (1971).
\newblock Dynamic limit pricing: Optimal pricing under threat of entry.
\newblock \emph{Journal of Economic Theory}, 3(3): 306--322.

\bibitem[Varian(1985)]{Varian1985}
Varian, H.R.\ (1985).
\newblock Price discrimination and social welfare.
\newblock \emph{American Economic Review}, 75(4): 870--875.

\end{thebibliography}

\end{document}
